\documentclass[aspectratio=169,10pt]{beamer}

% ============================================================
% CONFIGURACIÓN ROBUSTA
% ============================================================
\usepackage[utf8]{inputenc}
\usepackage[T1]{fontenc}
\usepackage[spanish,es-tabla]{babel}
\usepackage{lmodern}
\usepackage{microtype}
\usepackage{amsmath, amssymb, amsthm}
\usepackage{graphicx}
\usepackage{booktabs}
\usepackage{array}
\usepackage{multicol}
\usepackage{hyperref}
\usepackage{tikz}
\usetikzlibrary{shapes.geometric, arrows.meta, positioning, shadows, calc}
\usepackage{pifont}

% Gráficas numéricas
\usepackage{pgfplots}
\pgfplotsset{compat=1.18}
\usepgfplotslibrary{groupplots}

% Código
\usepackage{xcolor}
% Listings comentado debido a problemas de compatibilidad
% \usepackage{listings}

% ============================================================
% PALETA DE COLORES
% ============================================================
\definecolor{aerospace_blue}{RGB}{13,71,161}
\definecolor{aerospace_orange}{RGB}{230,81,0}
\definecolor{aerospace_cyan}{RGB}{0,151,167}
\definecolor{aerospace_gray}{RGB}{55,71,79}
\definecolor{code_bg}{RGB}{250,250,252}
\definecolor{code_keyword}{RGB}{13,71,161}
\definecolor{code_string}{RGB}{194,24,91}
\definecolor{code_comment}{RGB}{96,125,139}
\definecolor{highlight_yellow}{RGB}{255,243,224}

% ============================================================
% TEMA BEAMER
% ============================================================
\usetheme{Madrid}
\usecolortheme[named=aerospace_blue]{structure}

\setbeamercolor{palette primary}{bg=aerospace_blue,fg=white}
\setbeamercolor{palette secondary}{bg=aerospace_cyan,fg=white}
\setbeamercolor{palette tertiary}{bg=aerospace_orange,fg=white}
\setbeamercolor{frametitle}{bg=aerospace_blue,fg=white}
\setbeamercolor{block title}{bg=aerospace_cyan,fg=white}
\setbeamercolor{block body}{bg=highlight_yellow,fg=black}
\setbeamercolor{block title alerted}{bg=aerospace_orange,fg=white}
\setbeamercolor{block body alerted}{bg=aerospace_orange!10,fg=black}

\setbeamertemplate{navigation symbols}{}
\setbeamertemplate{page number in head/foot}[framenumber]

% ============================================================
% LOGO Y MARCA DE AGUA (CORREGIDO)
% ============================================================
\newcommand{\logofile}{logo_unii.png}

\setbeamertemplate{background}{
  \begin{tikzpicture}[remember picture,overlay]
    \node at (current page.center)
      [opacity=0.05] {\includegraphics[width=0.4\paperwidth]{\logofile}};
    \node[anchor=north east,opacity=0.7]
      at (current page.north east)
      {\includegraphics[height=0.8cm]{\logofile}};
  \end{tikzpicture}
}

% ============================================================
% LISTINGS (Python) - COMENTADO
% ============================================================
% \lstdefinestyle{pythonstyle}{
%   language=Python,
%   basicstyle=\ttfamily\footnotesize,
%   keywordstyle=\color{code_keyword}\bfseries,
%   stringstyle=\color{code_string},
%   commentstyle=\color{code_comment}\itshape,
%   backgroundcolor=\color{code_bg},
%   frame=single,
%   rulecolor=\color{aerospace_gray!30},
%   frameround=tttt,
%   breaklines=true,
%   showstringspaces=false,
%   numbers=left,
%   numberstyle=\tiny\color{aerospace_gray},
%   numbersep=8pt,
%   xleftmargin=15pt
% }

% ============================================================
% COMANDOS PERSONALIZADOS
% ============================================================
\newcommand{\importantbox}[2][Importante]{%
\begin{alertblock}{#1}#2\end{alertblock}}

\newcommand{\objectivebox}[1]{%
\begin{block}{Objetivo}#1\end{block}}

\newcommand{\dingcheck}{\textcolor{aerospace_cyan}{\ding{51}}}
\newcommand{\dingx}{\textcolor{aerospace_orange}{\ding{55}}}

% ============================================================
% FOOTER
% ============================================================
\setbeamertemplate{footline}{
  \leavevmode
  \hbox{
    \begin{beamercolorbox}[wd=.33\paperwidth,ht=2.5ex,dp=1ex,left]{author in head/foot}
      \hspace*{2ex}\insertshortauthor
    \end{beamercolorbox}
    \begin{beamercolorbox}[wd=.34\paperwidth,ht=2.5ex,dp=1ex,center]{title in head/foot}
      \insertshorttitle
    \end{beamercolorbox}
    \begin{beamercolorbox}[wd=.33\paperwidth,ht=2.5ex,dp=1ex,right]{date in head/foot}
      \insertframenumber{} / \inserttotalframenumber\hspace*{2ex}
    \end{beamercolorbox}
  }
}

% ============================================================
% DATOS DEL CURSO (CLASE 6)
% ============================================================
\title[Modelización Aeroespacial]{Modelización en Ingeniería Aeroespacial}
\subtitle{Clase 6: Discretización de Modelos y Estabilidad Numérica}
\author[Ing. Jair Molina]{Ing. Jair Molina Arce}
\institute[UNII]{Universidad Internacional de Innovación}
\date{\today}

% ============================================================
% TRANSICIÓN DE SECCIÓN
% ============================================================
\AtBeginSection[]{
\begin{frame}
\vfill
\centering
\begin{beamercolorbox}[sep=8pt,center,shadow=true,rounded=true]{title}
\usebeamerfont{title}\insertsectionhead
\end{beamercolorbox}
\vfill
\end{frame}
}

% ============================================================
% AJUSTES PGFPLOTS (estética neutra, legible)
% ============================================================
\pgfplotsset{
  every axis/.style={
    width=\linewidth,
    height=0.56\textheight,
    grid=both,
    major grid style={aerospace_gray!20},
    minor grid style={aerospace_gray!10},
    tick label style={font=\scriptsize},
    label style={font=\scriptsize},
    legend style={font=\scriptsize, draw=aerospace_gray!30, fill=white},
  }
}

% ============================================================
% DOCUMENTO
% ============================================================
\begin{document}

% ------------------------------------------------------------
\begin{frame}
\titlepage
\end{frame}

% ------------------------------------------------------------
\begin{frame}{Agenda}
\small
\begin{itemize}
\item ¿Por qué discretizar modelos continuos?
\item Modelos en tiempo continuo vs tiempo discreto
\item Discretización de sistemas lineales (exacta/ZOH)
\item Métodos numéricos como discretización (Euler, RK)
\item Estabilidad numérica y regiones de estabilidad
\item \textbf{Ejemplos con señales (continuo vs discreto)}
\item \textbf{Práctica: señal amortiguada (continua y discretizada)}
\item \textbf{Actividad: discretiza una señal propuesta y explora parámetros}
\end{itemize}
\end{frame}

% ============================================================
\section{Motivación}

\begin{frame}{¿Por qué discretizamos?}
\small

\objectivebox{
Entender cómo pasar de ecuaciones diferenciales a ecuaciones en diferencias
sin destruir la dinámica física del sistema.
}

\begin{block}{Realidad ingenieril}
\begin{itemize}
\item Computadoras trabajan en \textbf{tiempo discreto}
\item Sensores muestrean a intervalos $T_s$
\item Controladores digitales (MCU, FPGA, DSP)
\end{itemize}
\end{block}

\importantbox[Idea clave]{
Un sistema estable en continuo puede volverse \textbf{inestable} al discretizarlo mal.
}
\end{frame}

% ============================================================
\section{Modelos continuos vs discretos}

\begin{frame}{Modelo continuo}
\small
\begin{block}{Forma estándar}
\[
\dot{x}(t) = f(x(t),u(t))
\]
\end{block}

\begin{itemize}
\item Tiempo continuo $t \in \mathbb{R}$
\item Solución exacta rara vez disponible
\item Base física del sistema
\end{itemize}

\importantbox[Ejemplo aeroespacial]{
Dinámica de vuelo, actitud, órbitas, combustión, vibraciones.
}
\end{frame}

\begin{frame}{Modelo discreto}
\small
\begin{block}{Forma general}
\[
x_{k+1} = F(x_k, u_k)
\]
\end{block}

\begin{itemize}
\item Tiempo discreto $kT_s$
\item Aproximación del modelo continuo
\item Dependiente del método numérico
\end{itemize}

\importantbox[Advertencia]{
$F(\cdot)$ \textbf{no es único}: depende del método de discretización.
}
\end{frame}

% ============================================================
\section{Discretización de sistemas lineales}

\begin{frame}{Sistema lineal continuo}
\small
\begin{block}{Modelo}
\[
\dot{x} = A x + B u
\]
\end{block}

Queremos:
\[
x_{k+1} = A_d x_k + B_d u_k
\]

\importantbox[Pregunta central]{
¿Cómo obtener $A_d$, $B_d$ sin perder estabilidad?
}
\end{frame}

\begin{frame}{Discretización exacta (ZOH)}
\small
\begin{block}{Entrada constante por intervalo ($u(t)=u_k$)}
\[
A_d = e^{A T_s}
\]
\[
B_d = \int_0^{T_s} e^{A\tau} B\, d\tau
\]
\end{block}

\begin{itemize}
\item Preserva estabilidad y dinámica local
\item Es la forma correcta para control digital con ZOH
\end{itemize}

\importantbox[Regla]{
Si $\Re(\lambda_i(A))<0$ entonces típicamente $|\lambda_i(A_d)|<1$.
}
\end{frame}

% ============================================================
\section{Métodos numéricos como discretización}

\begin{frame}{Euler explícito}
\small
\begin{block}{Aproximación}
\[
x_{k+1} = x_k + T_s (A x_k + B u_k)
\]
\end{block}

\[
A_d = I + T_s A
\]

\begin{itemize}
\item Muy simple
\item Puede ser inestable para pasos grandes
\end{itemize}

\importantbox[Advertencia]{
Euler explícito puede volver inestable un sistema estable.
}
\end{frame}

\begin{frame}{Ejemplo: estable continuo $\rightarrow$ inestable discreto}
\small

\begin{block}{Sistema continuo}
\[
\dot{x}=-2x
\]
\end{block}

Continuo: $\lambda=-2$ (estable).\\[1mm]

\begin{block}{Euler}
\[
x_{k+1}=(1-2T_s)x_k
\]
\end{block}

\importantbox[Condición de estabilidad]{
\[
|1-2T_s|<1 \Rightarrow 0<T_s<1
\]
}

Si $T_s>1$, el método diverge aunque el sistema sea estable.
\end{frame}

% ============================================================
\section{Regiones de estabilidad}

\begin{frame}{Región de estabilidad: idea}
\small
\begin{block}{Criterio general}
Para el modo $\dot{x}=\lambda x$, el método es estable si:
\[
|R(z)|<1,\qquad z=\lambda T_s
\]
\end{block}

\begin{itemize}
\item $R(z)$: función de estabilidad del método
\item Estabilidad depende de $\lambda$ y del paso $T_s$
\end{itemize}

\importantbox[Lectura física]{
No basta con que $\lambda$ sea negativo: importa qué tanto ``cabe'' en el paso.
}
\end{frame}

\begin{frame}{Comparación práctica de métodos}
\small
\begin{block}{Resumen (intuición)}
\begin{itemize}
\item \textbf{Euler explícito:} región pequeña (peligroso en modos rápidos)
\item \textbf{RK4:} región más grande (robusto para muchos sistemas)
\item \textbf{Implícitos:} región grande (ideal para sistemas stiff)
\end{itemize}
\end{block}

\importantbox[Regla de ingeniería]{
Si el sistema tiene modos rápidos:
\begin{itemize}
\item reduce $T_s$ o
\item cambia el integrador
\end{itemize}
}
\end{frame}

% ============================================================
\section{Ejemplos con señales}

% ------------------ Ejemplo 1 (Subframe 1/2) ------------------
\begin{frame}{Ejemplo 1 (1/2): seno continuo vs muestreado}
\small

\begin{block}{Señal continua}
\[
x(t)=A\sin(2\pi f t)
\]
\end{block}

\begin{block}{Señal discretizada (muestreo)}
\[
x[k]=x(kT_s)=A\sin\big(2\pi f\,kT_s\big)
\]
\end{block}

\begin{itemize}
\item Si $f$ es alta y $T_s$ es grande, aparecen distorsiones (aliasing).
\item En control digital y sensores: el muestreo define lo que \textbf{ves} del fenómeno.
\end{itemize}

\importantbox[Tip]{
Usa $f_s=\frac{1}{T_s}$ suficientemente alto respecto a $f$ para representar bien la señal.
}
\end{frame}

% ------------------ Ejemplo 1 (Subframe 2/2) ------------------
\begin{frame}{Ejemplo 1 (2/2): gráfica seno continuo vs muestreado}
\small
\begin{columns}[T]
\column{0.34\textwidth}
\begin{block}{Parámetros (demo)}
\begin{itemize}
\item $A=1$
\item $f=2\,$Hz
\item $T_s=0.15\,$s
\item $t \in [0,2]\,$s
\end{itemize}
\end{block}
\importantbox[Lectura]{
Marcadores separados $\Rightarrow$ pocos puntos por ciclo.
}

\column{0.66\textwidth}
\begin{tikzpicture}
\begin{axis}[
  xlabel={$t$ [s]},
  ylabel={$x$},
  ymin=-1.2, ymax=1.2,
  legend pos=south east,
]
% Continuo
\addplot+[thick, samples=500, domain=0:2]
  {sin(deg(2*pi*2*x))};
\addlegendentry{Continuo $x(t)$}

% Discreto (coordenadas)
\def\Ts{0.15}
\addplot+[ycomb, mark=*, samples=1]
  coordinates {
    (0.00,{sin(deg(2*pi*2*0.00))})
    (0.15,{sin(deg(2*pi*2*0.15))})
    (0.30,{sin(deg(2*pi*2*0.30))})
    (0.45,{sin(deg(2*pi*2*0.45))})
    (0.60,{sin(deg(2*pi*2*0.60))})
    (0.75,{sin(deg(2*pi*2*0.75))})
    (0.90,{sin(deg(2*pi*2*0.90))})
    (1.05,{sin(deg(2*pi*2*1.05))})
    (1.20,{sin(deg(2*pi*2*1.20))})
    (1.35,{sin(deg(2*pi*2*1.35))})
    (1.50,{sin(deg(2*pi*2*1.50))})
    (1.65,{sin(deg(2*pi*2*1.65))})
    (1.80,{sin(deg(2*pi*2*1.80))})
    (1.95,{sin(deg(2*pi*2*1.95))})
  };
\addlegendentry{Discreto $x[k]$}
\end{axis}
\end{tikzpicture}
\end{columns}
\end{frame}

% ------------------ Ejemplo 2 (Subframe 1/2) ------------------
\begin{frame}{Ejemplo 2 (1/2): exponencial (decaimiento) continuo vs discreto}
\small

\begin{block}{Señal continua}
\[
x(t)=A e^{-a t},\qquad a>0
\]
\end{block}

\begin{block}{Muestreo directo}
\[
x[k]=A e^{-a k T_s}
\]
\end{block}

\begin{itemize}
\item Es un caso muy limpio: el muestreo \textbf{no cambia} la estabilidad (solo la resolución).
\item Sirve como referencia para entender señales amortiguadas.
\end{itemize}

\importantbox[Idea]{
Aquí no hay oscilación: solo una escala de tiempo $\tau=\frac{1}{a}$.
}
\end{frame}

% ------------------ Ejemplo 2 (Subframe 2/2) ------------------
\begin{frame}{Ejemplo 2 (2/2): gráfica exponencial continuo vs muestreado}
\small
\begin{columns}[T]
\column{0.34\textwidth}
\begin{block}{Parámetros (demo)}
\begin{itemize}
\item $A=1$
\item $a=1.5\,$s$^{-1}$
\item $T_s=0.25\,$s
\item $t \in [0,3]\,$s
\end{itemize}
\end{block}
\importantbox[Lectura]{
La forma se conserva; cambian los puntos disponibles.
}

\column{0.66\textwidth}
\begin{tikzpicture}
\begin{axis}[
  xlabel={$t$ [s]},
  ylabel={$x$},
  ymin=0, ymax=1.05,
  legend pos=north east,
]
% Continuo
\addplot+[thick, samples=500, domain=0:3]
  {exp(-1.5*x)};
\addlegendentry{Continuo $x(t)$}

% Discreto
\addplot+[ycomb, mark=*, samples=1]
  coordinates {
    (0.00,{exp(-1.5*0.00)})
    (0.25,{exp(-1.5*0.25)})
    (0.50,{exp(-1.5*0.50)})
    (0.75,{exp(-1.5*0.75)})
    (1.00,{exp(-1.5*1.00)})
    (1.25,{exp(-1.5*1.25)})
    (1.50,{exp(-1.5*1.50)})
    (1.75,{exp(-1.5*1.75)})
    (2.00,{exp(-1.5*2.00)})
    (2.25,{exp(-1.5*2.25)})
    (2.50,{exp(-1.5*2.50)})
    (2.75,{exp(-1.5*2.75)})
    (3.00,{exp(-1.5*3.00)})
  };
\addlegendentry{Discreto $x[k]$}
\end{axis}
\end{tikzpicture}
\end{columns}
\end{frame}

% ============================================================
\section{Práctica: señal amortiguada (continua y discretizada)}

% ------------------ Amortiguada (Subframe 1/2) ------------------
\begin{frame}{Práctica (1/2): señal amortiguada (modelo continuo y muestreo)}
\small

\objectivebox{
Comparar una señal amortiguada en continuo vs su versión discretizada.
}

\begin{block}{Señal (oscilación amortiguada)}
\[
x(t)=A\,e^{-\alpha t}\,\sin(2\pi f t+\varphi)
\]
\end{block}

\begin{block}{Muestreo ideal}
\[
x[k]=x(kT_s)=A\,e^{-\alpha kT_s}\,\sin(2\pi f\,kT_s+\varphi)
\]
\end{block}

\begin{itemize}
\item $A$: amplitud inicial
\item $\alpha$: amortiguamiento (1/s)
\item $f$: frecuencia (Hz)
\item $\varphi$: fase (rad)
\end{itemize}

\importantbox[Lectura física]{
La envolvente $A e^{-\alpha t}$ impone la rapidez con la que ``muere'' la oscilación.
}
\end{frame}

% ------------------ Amortiguada (Subframe 2/2) ------------------
\begin{frame}{Práctica (2/2): gráfica señal amortiguada continuo vs muestreado}
\small
\begin{columns}[T]
\column{0.34\textwidth}
\begin{block}{Parámetros (demo)}
\begin{itemize}
\item $A=1$
\item $\alpha=0.7\,$s$^{-1}$
\item $f=2\,$Hz
\item $\varphi=0$
\item $T_s=0.12\,$s
\item $t \in [0,5]\,$s
\end{itemize}
\end{block}
\importantbox[Objetivo]{
Elegir $T_s$ para conservar forma, fase y envolvente.
}

\column{0.66\textwidth}
\begin{tikzpicture}
\begin{axis}[
  xlabel={$t$ [s]},
  ylabel={$x$},
  ymin=-1.1, ymax=1.1,
  legend pos=south east,
]
% Continuo
\addplot+[thick, samples=600, domain=0:5]
  {exp(-0.7*x)*sin(deg(2*pi*2*x))};
\addlegendentry{Continuo $x(t)$}

% Discreto
\addplot+[ycomb, mark=*, samples=1]
  coordinates {
    (0.00,{exp(-0.7*0.00)*sin(deg(2*pi*2*0.00))})
    (0.12,{exp(-0.7*0.12)*sin(deg(2*pi*2*0.12))})
    (0.24,{exp(-0.7*0.24)*sin(deg(2*pi*2*0.24))})
    (0.36,{exp(-0.7*0.36)*sin(deg(2*pi*2*0.36))})
    (0.48,{exp(-0.7*0.48)*sin(deg(2*pi*2*0.48))})
    (0.60,{exp(-0.7*0.60)*sin(deg(2*pi*2*0.60))})
    (0.72,{exp(-0.7*0.72)*sin(deg(2*pi*2*0.72))})
    (0.84,{exp(-0.7*0.84)*sin(deg(2*pi*2*0.84))})
    (0.96,{exp(-0.7*0.96)*sin(deg(2*pi*2*0.96))})
    (1.08,{exp(-0.7*1.08)*sin(deg(2*pi*2*1.08))})
    (1.20,{exp(-0.7*1.20)*sin(deg(2*pi*2*1.20))})
    (1.32,{exp(-0.7*1.32)*sin(deg(2*pi*2*1.32))})
    (1.44,{exp(-0.7*1.44)*sin(deg(2*pi*2*1.44))})
    (1.56,{exp(-0.7*1.56)*sin(deg(2*pi*2*1.56))})
    (1.68,{exp(-0.7*1.68)*sin(deg(2*pi*2*1.68))})
    (1.80,{exp(-0.7*1.80)*sin(deg(2*pi*2*1.80))})
    (1.92,{exp(-0.7*1.92)*sin(deg(2*pi*2*1.92))})
    (2.04,{exp(-0.7*2.04)*sin(deg(2*pi*2*2.04))})
    (2.16,{exp(-0.7*2.16)*sin(deg(2*pi*2*2.16))})
    (2.28,{exp(-0.7*2.28)*sin(deg(2*pi*2*2.28))})
    (2.40,{exp(-0.7*2.40)*sin(deg(2*pi*2*2.40))})
    (2.52,{exp(-0.7*2.52)*sin(deg(2*pi*2*2.52))})
    (2.64,{exp(-0.7*2.64)*sin(deg(2*pi*2*2.64))})
    (2.76,{exp(-0.7*2.76)*sin(deg(2*pi*2*2.76))})
    (2.88,{exp(-0.7*2.88)*sin(deg(2*pi*2*2.88))})
    (3.00,{exp(-0.7*3.00)*sin(deg(2*pi*2*3.00))})
    (3.12,{exp(-0.7*3.12)*sin(deg(2*pi*2*3.12))})
    (3.24,{exp(-0.7*3.24)*sin(deg(2*pi*2*3.24))})
    (3.36,{exp(-0.7*3.36)*sin(deg(2*pi*2*3.36))})
    (3.48,{exp(-0.7*3.48)*sin(deg(2*pi*2*3.48))})
    (3.60,{exp(-0.7*3.60)*sin(deg(2*pi*2*3.60))})
    (3.72,{exp(-0.7*3.72)*sin(deg(2*pi*2*3.72))})
    (3.84,{exp(-0.7*3.84)*sin(deg(2*pi*2*3.84))})
    (3.96,{exp(-0.7*3.96)*sin(deg(2*pi*2*3.96))})
    (4.08,{exp(-0.7*4.08)*sin(deg(2*pi*2*4.08))})
    (4.20,{exp(-0.7*4.20)*sin(deg(2*pi*2*4.20))})
    (4.32,{exp(-0.7*4.32)*sin(deg(2*pi*2*4.32))})
    (4.44,{exp(-0.7*4.44)*sin(deg(2*pi*2*4.44))})
    (4.56,{exp(-0.7*4.56)*sin(deg(2*pi*2*4.56))})
    (4.68,{exp(-0.7*4.68)*sin(deg(2*pi*2*4.68))})
    (4.80,{exp(-0.7*4.80)*sin(deg(2*pi*2*4.80))})
    (4.92,{exp(-0.7*4.92)*sin(deg(2*pi*2*4.92))})
  };
\addlegendentry{Discreto $x[k]$}
\end{axis}
\end{tikzpicture}
\end{columns}
\end{frame}

% ============================================================
\section{Código Python de la práctica}

\begin{frame}{Estructura del script (visión general)}
\small
\begin{block}{Bloques del programa}
\begin{enumerate}
\item Parámetros (A, alpha, f, phi, Ts, duración)
\item Generación de malla continua $t$ (alta resolución)
\item Generación de instantes discretos $t_k=kT_s$
\item Cálculo $x(t)$ y $x[k]$
\item Gráficas: continuo (línea) + discreto (stems o marcadores)
\end{enumerate}
\end{block}

\importantbox[Tip de ingeniería]{
Mismo modelo físico, dos representaciones: lo que cambia es el \textbf{tiempo} (continuo vs muestreado).
}
\end{frame}

\begin{frame}[fragile]{Código Python: señales y muestreo}
\small
\begin{block}{Funciones principales}
\texttt{damped\_signal(t, A, alpha, f, phi)}\\
\quad Calcula: $x(t) = A \cdot e^{-\alpha t} \cdot \sin(2\pi f t + \varphi)$

\vspace{2mm}
\texttt{sample\_times(Ts, t0, tf)}\\
\quad Genera instantes discretos: $t_k = t_0 + k \cdot T_s$

\vspace{2mm}
\texttt{main()}\\
\quad Configura parámetros, calcula señales y grafica
\end{block}

\importantbox[Código disponible]{
El código Python completo está disponible en el repositorio del curso. Incluye NumPy, Matplotlib y ejemplos completos para experimentar con diferentes parámetros.
}
\end{frame}

% ============================================================
\section{Actividad}

% ------------------ Actividad chirp (Subframe 1/2) ------------------
\begin{frame}{Actividad (1/2): señal propuesta para discretizar}
\small

\objectivebox{
Cada alumno debe discretizar la señal, graficar continuo vs discreto y justificar el $T_s$ elegido.
}

\begin{block}{Señal propuesta (chirp lineal amortiguado)}
\[
x(t)=A\,e^{-\alpha t}\,\sin\!\Big(2\pi\big(f_0 t + \tfrac{1}{2}\beta t^2\big)+\varphi\Big)
\]
\end{block}

\begin{itemize}
\item $f_0$: frecuencia inicial (Hz)
\item $\beta$: tasa de barrido (Hz/s)
\item Misma idea: $x[k]=x(kT_s)$
\end{itemize}

\importantbox[Qué deben modificar]{
$A,\alpha,f_0,\beta,\varphi,T_s$ y observar aliasing vs fidelidad.
}
\end{frame}

% ------------------ Actividad chirp (Subframe 2/2) ------------------
\begin{frame}{Actividad (2/2): gráfica ejemplo (chirp amortiguado) continuo vs muestreado}
\small
\begin{columns}[T]
\column{0.34\textwidth}
\begin{block}{Parámetros (demo)}
\begin{itemize}
\item $A=1$
\item $\alpha=0.25\,$s$^{-1}$
\item $f_0=1\,$Hz
\item $\beta=2.0\,$Hz/s
\item $\varphi=0$
\item $T_s=0.08\,$s
\item $t\in[0,5]\,$s
\end{itemize}
\end{block}
\importantbox[Lectura]{
A mayor $t$, sube $f(t)=f_0+\beta t$; si $f_s$ es bajo, aparece aliasing.
}

\column{0.66\textwidth}
\begin{tikzpicture}
\begin{axis}[
  xlabel={$t$ [s]},
  ylabel={$x$},
  ymin=-1.1, ymax=1.1,
  legend pos=south east,
]
% Continuo: x(t)=exp(-alpha t)*sin(2pi(f0 t + 0.5 beta t^2))
\addplot+[thick, samples=700, domain=0:5]
  {exp(-0.25*x)*sin(deg(2*pi*(1*x + 0.5*2.0*x^2)))};
\addlegendentry{Continuo $x(t)$}

% Discreto (coordenadas)
\addplot+[ycomb, mark=*, samples=1]
  coordinates {
    (0.00,{exp(-0.25*0.00)*sin(deg(2*pi*(1*0.00 + 0.5*2.0*0.00^2)))})
    (0.08,{exp(-0.25*0.08)*sin(deg(2*pi*(1*0.08 + 0.5*2.0*0.08^2)))})
    (0.16,{exp(-0.25*0.16)*sin(deg(2*pi*(1*0.16 + 0.5*2.0*0.16^2)))})
    (0.24,{exp(-0.25*0.24)*sin(deg(2*pi*(1*0.24 + 0.5*2.0*0.24^2)))})
    (0.32,{exp(-0.25*0.32)*sin(deg(2*pi*(1*0.32 + 0.5*2.0*0.32^2)))})
    (0.40,{exp(-0.25*0.40)*sin(deg(2*pi*(1*0.40 + 0.5*2.0*0.40^2)))})
    (0.48,{exp(-0.25*0.48)*sin(deg(2*pi*(1*0.48 + 0.5*2.0*0.48^2)))})
    (0.56,{exp(-0.25*0.56)*sin(deg(2*pi*(1*0.56 + 0.5*2.0*0.56^2)))})
    (0.64,{exp(-0.25*0.64)*sin(deg(2*pi*(1*0.64 + 0.5*2.0*0.64^2)))})
    (0.72,{exp(-0.25*0.72)*sin(deg(2*pi*(1*0.72 + 0.5*2.0*0.72^2)))})
    (0.80,{exp(-0.25*0.80)*sin(deg(2*pi*(1*0.80 + 0.5*2.0*0.80^2)))})
    (0.88,{exp(-0.25*0.88)*sin(deg(2*pi*(1*0.88 + 0.5*2.0*0.88^2)))})
    (0.96,{exp(-0.25*0.96)*sin(deg(2*pi*(1*0.96 + 0.5*2.0*0.96^2)))})
    (1.04,{exp(-0.25*1.04)*sin(deg(2*pi*(1*1.04 + 0.5*2.0*1.04^2)))})
    (1.12,{exp(-0.25*1.12)*sin(deg(2*pi*(1*1.12 + 0.5*2.0*1.12^2)))})
    (1.20,{exp(-0.25*1.20)*sin(deg(2*pi*(1*1.20 + 0.5*2.0*1.20^2)))})
    (1.28,{exp(-0.25*1.28)*sin(deg(2*pi*(1*1.28 + 0.5*2.0*1.28^2)))})
    (1.36,{exp(-0.25*1.36)*sin(deg(2*pi*(1*1.36 + 0.5*2.0*1.36^2)))})
    (1.44,{exp(-0.25*1.44)*sin(deg(2*pi*(1*1.44 + 0.5*2.0*1.44^2)))})
    (1.52,{exp(-0.25*1.52)*sin(deg(2*pi*(1*1.52 + 0.5*2.0*1.52^2)))})
    (1.60,{exp(-0.25*1.60)*sin(deg(2*pi*(1*1.60 + 0.5*2.0*1.60^2)))})
    (1.68,{exp(-0.25*1.68)*sin(deg(2*pi*(1*1.68 + 0.5*2.0*1.68^2)))})
    (1.76,{exp(-0.25*1.76)*sin(deg(2*pi*(1*1.76 + 0.5*2.0*1.76^2)))})
    (1.84,{exp(-0.25*1.84)*sin(deg(2*pi*(1*1.84 + 0.5*2.0*1.84^2)))})
    (1.92,{exp(-0.25*1.92)*sin(deg(2*pi*(1*1.92 + 0.5*2.0*1.92^2)))})
    (2.00,{exp(-0.25*2.00)*sin(deg(2*pi*(1*2.00 + 0.5*2.0*2.00^2)))})
    (2.08,{exp(-0.25*2.08)*sin(deg(2*pi*(1*2.08 + 0.5*2.0*2.08^2)))})
    (2.16,{exp(-0.25*2.16)*sin(deg(2*pi*(1*2.16 + 0.5*2.0*2.16^2)))})
    (2.24,{exp(-0.25*2.24)*sin(deg(2*pi*(1*2.24 + 0.5*2.0*2.24^2)))})
    (2.32,{exp(-0.25*2.32)*sin(deg(2*pi*(1*2.32 + 0.5*2.0*2.32^2)))})
    (2.40,{exp(-0.25*2.40)*sin(deg(2*pi*(1*2.40 + 0.5*2.0*2.40^2)))})
    (2.48,{exp(-0.25*2.48)*sin(deg(2*pi*(1*2.48 + 0.5*2.0*2.48^2)))})
    (2.56,{exp(-0.25*2.56)*sin(deg(2*pi*(1*2.56 + 0.5*2.0*2.56^2)))})
    (2.64,{exp(-0.25*2.64)*sin(deg(2*pi*(1*2.64 + 0.5*2.0*2.64^2)))})
    (2.72,{exp(-0.25*2.72)*sin(deg(2*pi*(1*2.72 + 0.5*2.0*2.72^2)))})
    (2.80,{exp(-0.25*2.80)*sin(deg(2*pi*(1*2.80 + 0.5*2.0*2.80^2)))})
    (2.88,{exp(-0.25*2.88)*sin(deg(2*pi*(1*2.88 + 0.5*2.0*2.88^2)))})
    (2.96,{exp(-0.25*2.96)*sin(deg(2*pi*(1*2.96 + 0.5*2.0*2.96^2)))})
    (3.04,{exp(-0.25*3.04)*sin(deg(2*pi*(1*3.04 + 0.5*2.0*3.04^2)))})
    (3.12,{exp(-0.25*3.12)*sin(deg(2*pi*(1*3.12 + 0.5*2.0*3.12^2)))})
    (3.20,{exp(-0.25*3.20)*sin(deg(2*pi*(1*3.20 + 0.5*2.0*3.20^2)))})
    (3.28,{exp(-0.25*3.28)*sin(deg(2*pi*(1*3.28 + 0.5*2.0*3.28^2)))})
    (3.36,{exp(-0.25*3.36)*sin(deg(2*pi*(1*3.36 + 0.5*2.0*3.36^2)))})
    (3.44,{exp(-0.25*3.44)*sin(deg(2*pi*(1*3.44 + 0.5*2.0*3.44^2)))})
    (3.52,{exp(-0.25*3.52)*sin(deg(2*pi*(1*3.52 + 0.5*2.0*3.52^2)))})
    (3.60,{exp(-0.25*3.60)*sin(deg(2*pi*(1*3.60 + 0.5*2.0*3.60^2)))})
    (3.68,{exp(-0.25*3.68)*sin(deg(2*pi*(1*3.68 + 0.5*2.0*3.68^2)))})
    (3.76,{exp(-0.25*3.76)*sin(deg(2*pi*(1*3.76 + 0.5*2.0*3.76^2)))})
    (3.84,{exp(-0.25*3.84)*sin(deg(2*pi*(1*3.84 + 0.5*2.0*3.84^2)))})
    (3.92,{exp(-0.25*3.92)*sin(deg(2*pi*(1*3.92 + 0.5*2.0*3.92^2)))})
    (4.00,{exp(-0.25*4.00)*sin(deg(2*pi*(1*4.00 + 0.5*2.0*4.00^2)))})
    (4.08,{exp(-0.25*4.08)*sin(deg(2*pi*(1*4.08 + 0.5*2.0*4.08^2)))})
    (4.16,{exp(-0.25*4.16)*sin(deg(2*pi*(1*4.16 + 0.5*2.0*4.16^2)))})
    (4.24,{exp(-0.25*4.24)*sin(deg(2*pi*(1*4.24 + 0.5*2.0*4.24^2)))})
    (4.32,{exp(-0.25*4.32)*sin(deg(2*pi*(1*4.32 + 0.5*2.0*4.32^2)))})
    (4.40,{exp(-0.25*4.40)*sin(deg(2*pi*(1*4.40 + 0.5*2.0*4.40^2)))})
    (4.48,{exp(-0.25*4.48)*sin(deg(2*pi*(1*4.48 + 0.5*2.0*4.48^2)))})
    (4.56,{exp(-0.25*4.56)*sin(deg(2*pi*(1*4.56 + 0.5*2.0*4.56^2)))})
    (4.64,{exp(-0.25*4.64)*sin(deg(2*pi*(1*4.64 + 0.5*2.0*4.64^2)))})
    (4.72,{exp(-0.25*4.72)*sin(deg(2*pi*(1*4.72 + 0.5*2.0*4.72^2)))})
    (4.80,{exp(-0.25*4.80)*sin(deg(2*pi*(1*4.80 + 0.5*2.0*4.80^2)))})
    (4.88,{exp(-0.25*4.88)*sin(deg(2*pi*(1*4.88 + 0.5*2.0*4.88^2)))})
    (4.96,{exp(-0.25*4.96)*sin(deg(2*pi*(1*4.96 + 0.5*2.0*4.96^2)))})
  };
\addlegendentry{Discreto $x[k]$}
\end{axis}
\end{tikzpicture}
\end{columns}
\end{frame}

% ============================================================
\section{Sistemas rígidos (Stiff systems)}

\begin{frame}{¿Qué es un sistema stiff?}
\small
\begin{block}{Definición práctica}
Sistema con escalas de tiempo muy distintas:
\[
|\Re(\lambda_{\text{rápido}})| \gg |\Re(\lambda_{\text{lento}})|
\]
\end{block}

\begin{itemize}
\item La estabilidad obliga a pasos muy pequeños
\item La precisión no es el problema: la estabilidad numérica sí
\end{itemize}

\importantbox[Ejemplos aeroespaciales]{
\begin{itemize}
\item combustión (rápida) + estructura (lenta)
\item térmico (lento) + actuadores (rápidos)
\item vibración local (rápida) + trayectoria (lenta)
\end{itemize}
}
\end{frame}

\begin{frame}{Qué se observa en la práctica}
\small
\begin{block}{Síntomas típicos}
\begin{itemize}
\item La simulación ``explota'' con Euler, aunque la física sea estable
\item RK4 requiere pasos muy pequeños para no divergir
\item El tiempo de cómputo se dispara
\end{itemize}
\end{block}

\importantbox[Soluciones]{
\begin{itemize}
\item Integradores implícitos (Backward Euler, BDF)
\item RK adaptativos (solve\_ivp con métodos stiff: Radau, BDF)
\item Separación de escalas / reducción de modelo
\end{itemize}
}
\end{frame}

% ============================================================
\section{Consecuencias en control digital}

\begin{frame}{Control en tiempo discreto}
\small
\begin{block}{Realidad de implementación}
\begin{itemize}
\item Control diseñado en continuo
\item Implementado en discreto con $T_s$
\item ZOH y discretización afectan polos
\end{itemize}
\end{block}

\importantbox[Regla empírica]{
Para capturar dinámica sin distorsión severa:
\[
T_s \le \frac{1}{10}\, \tau_{\text{modo más rápido}}
\]
}
\end{frame}

\begin{frame}{Mensaje final: simulación vs implementación}
\small
\begin{columns}[T]
\column{0.52\textwidth}
\begin{block}{Simulación}
\begin{itemize}
\item objetivo: entender física
\item tolera pasos pequeños
\item comparas integradores
\end{itemize}
\end{block}

\column{0.48\textwidth}
\begin{block}{Implementación (MCU)}
\begin{itemize}
\item objetivo: ejecutar rápido y estable
\item $T_s$ limitado por hardware
\item discretización correcta es crítica
\end{itemize}
\end{block}
\end{columns}

\importantbox[Regla]{
Si el resultado cambia drásticamente al cambiar $T_s$, aún no dominas el modelo.
}
\end{frame}

% ============================================================
\section{Resumen}

\begin{frame}{Lecciones clave de la Clase 6}
\small
\begin{enumerate}
\item Discretizar es una decisión física y numérica
\item Euler explícito puede destruir estabilidad
\item La estabilidad numérica depende de $z=\lambda T_s$
\item Los sistemas stiff requieren métodos adecuados
\item Muestreo $T_s$ define fidelidad de señales y control digital
\end{enumerate}

\vspace{3mm}
\begin{center}
\Large \textcolor{aerospace_blue}{\textbf{Siguiente clase:}}\\
\large Control lineal: feedback, polos y diseño (introducción)
\end{center}
\end{frame}

% ============================================================
\section*{Preguntas}

\begin{frame}
\vfill
\centering
\begin{beamercolorbox}[sep=8pt,center,shadow=true,rounded=true]{title}
\usebeamerfont{title}\Huge ¿Preguntas?
\end{beamercolorbox}
\vfill

\begin{center}
\large
\texttt{jair.molina@unii.edu.mx}\\[2mm]
\small
Material disponible en el repositorio del curso
\end{center}
\end{frame}

\end{document}
